\chapter{Appendix}

\section{Coding Framework Overview}
This appendix provides a summary of the thematic coding framework used in the dissertation. The coding process was informed by Puntoni et al.'s experiential framework and translated into four overarching consumer-perception tensions. The appendix is included to show how the main analytical categories used in Chapter 3 and Chapter 4 were derived from the literature.

\begin{table}[h!]
\centering
\caption{Summary of the thematic coding framework}
\begin{tabular}{p{3cm}p{3.2cm}p{3.3cm}p{4cm}}
\toprule
\textbf{Theme group} & \textbf{Puntoni domain} & \textbf{Core tension} & \textbf{Illustrative constructs} \\
\midrule
Data capture and privacy & Data capture experience & Served versus exploited & Privacy calculus, dataveillance, perceived vulnerability, transparency, trust \\
\midrule
Classification and personalisation & Classification experience & Understood versus misunderstood & Personalisation, uniqueness neglect, expected group size, stereotyping, algorithmic fairness \\
\midrule
Task delegation and autonomy & Delegation experience & Empowered versus replaced & Self-efficacy, perceived control, perceived sacrifice, autonomy, algorithm aversion \\
\midrule
Social interaction and authenticity & Social experience & Connected versus alienated & Source authenticity, emotional authenticity, empathy gap, moral disgust, novelty \\
\bottomrule
\end{tabular}
\end{table}

\section{Coding Table}
Table A1 provides a fuller version of the coding logic used in the dissertation. It links the main theme groups to associated concepts, supporting studies, and their interpretive relevance to consumer perception.

\begin{table}[h!]
\centering
\caption{Full thematic coding table}
\begin{tabular}{p{2.8cm}p{3.4cm}p{4.4cm}p{4cm}}
\toprule
\textbf{Theme group} & \textbf{Key constructs} & \textbf{Illustrative supporting studies} & \textbf{Interpretive relevance} \\
\midrule
Data capture and privacy & Privacy calculus, surveillance, dataveillance, trust, transparency, perceived exploitation & Aguirre et al. (2015); Boerman et al. (2017); Martin and Murphy (2017); Hardcastle et al. (2025); An and Ngo (2025); Srivastava and Gurme (2026); Ozturkcan and Bozda{\u{g}} (2025) & Shows how consumers may value tailored service while simultaneously resisting opaque or excessive data capture \\
\midrule
Classification and personalisation & Personalisation, algorithmic classification, uniqueness neglect, expected group size, bias, fairness & Puntoni et al. (2021); Hou et al. (2026); Starke et al. (2022) & Shows that AI may create relevance without necessarily creating a feeling of individuality or fair recognition \\
\midrule
Task delegation and autonomy & Delegation, self-efficacy, autonomy, perceived control, replacement, psychological threat & Puntoni et al. (2021); Davenport et al. (2020); Ameen et al. (2021); Hardcastle et al. (2025) & Shows how AI can reduce effort while also narrowing choice and weakening the consumer's sense of agency \\
\midrule
Social interaction and authenticity & Source authenticity, emotional authenticity, empathy gap, novelty, alienation, moral disgust & Kirk and Givi (2025); Shi and Jiang (2026); Arango et al. (2023); Belanche et al. (2025); Longoni et al. (2019) & Shows how AI-mediated communication may be accepted in factual contexts but resisted in emotional or high-involvement settings \\
\bottomrule
\end{tabular}
\end{table}

\section{Appendix Note}
The appendix coding tables are included to make the analytical process transparent. They are not intended to replace the discussion in the main chapters, but to document how the dissertation moved from literature selection to thematic interpretation.
